\section{Related Work}
\label{sec:related}

\paragraph{Conflict-handling styles.} The dual-concern tradition describes conflict behaviour by concern for one's own and for the other's outcomes: the managerial grid \citep{blake1964managerial}, the \tki{} model and instrument \citep{thomas1974tki,thomas1992conflict}, Rahim's inventory \citep{rahim1983measure} and the dual-concern model of negotiation \citep{pruitt1986social}; \citet{vandevliert1990toward} assess the instruments against the theory. We test the model's situational prescriptions \citep{thomas1977toward} where a strategy can be assigned rather than reported.

\paragraph{LLM agents in negotiation.} Negotiation dialogue was studied before LLMs with learned agents \citep{lewis2017deal,he2018decoupling} and human corpora such as CaSiNo \citep{chawla2021casino}. Recent work evaluates LLMs as negotiators in arenas and benchmarks \citep{bianchi2024negotiationarena,abdelnabi2024cooperation,xia2024measuring,kwon2024effective,davidson2024evaluating}, improves them by self-play and feedback \citep{fu2023improving}, studies repeated games and strategic reasoning \citep{akata2023playing,gandhi2023strategic} and personas in negotiation prompts \citep{noh2024llms}; \citet{zhan2024negotiate} survey the area. These ask how well a model negotiates or which prompt wins a given game. We ask which conflict-handling strategy to assign as a function of the situation, with compliance checked from the moves and with efficiency and cost on one scale, and we reuse two of these benchmarks as external settings.

\paragraph{Game-theoretic workflows.} \citet{hua2024gametheoretic} ask whether LLMs are rational in matrix games and in Deal-or-No-Deal and propose workflows that make them so. We manipulate a different variable (an assigned handling strategy, not a reasoning workflow), measure a different outcome (the distribution of payoff, cost and efficiency over strategies and conditions, not distance from equilibrium) and propose no method. Two of their findings are starting points here: their ``undue trust'' is our default agent leaving a tenth of the joint value on the table, and their one-sided workflow being exploited is one cell of our matrix.

\paragraph{Coordination between coding agents.} CooperBench \citep{cooperbench} and ParallACT \citep{parallact} study agents that edit the same code; conflicts there are binary (whose change lands), so conceding is the only cooperative move. Our grounded layer reduces ParallACT's contradictory ticket pairs to the design decisions behind them, where there is something to trade.

